\chapter{Abstandssumme gegen Abstandsdifferenz}
Nun wird die Verteilung der berechneten Abstände zu den Anodendrähten betrachtet. Hierzu wird für Ereignisse, bei denen zwei benachbarte Zellen angesprochen haben, die Summen der Abstände gegen die halben Differenzen der Abstände aufgetragen. Die Abstände werden hierbei über die zuvor bestimmte ODB berechnet. Diese betrachtung wird für verschiedene Winkelbereiche durchgeführt. Für Winkelbereiche nahe der Flächennormalen, wie in Abbildung \ref{fig:sum_diff_center} und \ref{fig:sum_diff_mid} erkennt man eine Häufung bei entlang der X-Achse bei einer Abstandssumme von etwa $\SI{8}{\milli\meter}$ bis etwa $\SI{12}{\milli\meter}$. Dies entspricht der Erwartung, dass die Summe der Abstände bei normalem Einfall der dem Drahtabstand, also $\SI{8.5}{\milli\meter}$ entsprechen muss. Bei Flacheren Einfallwinkeln erwartet man, dass die Summe der Abstände größer wird, da die Zellen schräg durchflogen werden. Dies findet sich in den Abbildungen \ref{fig:sum_diff_all} und \ref{fig:sum_diff_edge} in Form einer Anhäufung bei einer Differenz von ca. Null und einer Summe von etwa $\SI{14}{\milli\meter}$ wieder. Je extremer die betrachteten Winkel werden, desto deutlicher lassen sich die jeweiligen Strukturen erkennen, bei gleichzeitig weniger gesamtereignissen. Beides entspricht den Erwartungen.

\jafpp{sum_diff_all}{Auftragung Abstandssumme gegen halbe Abstandsdifferenz für alle Winkel}{sum_diff_mid}{ca. $\SI{-18}{\degree}$ bis ca. $\SI{14}{\degree}}
\jafpp{sum_diff_center}{Auftragung Abstandssumme gegen halbe Abstandsdifferenz für ca. $\SI{-9}{\degree}$ bis ca. $\SI{5}{\degree}$}{sum_diff_edge}{für ca. $\SI{42}{\degree}$ bis ca. $\SI{64}{\degree}$}
